\chapter{WIS - Weighted Interval Scheduling 
\\(ottimizzazione combinatoria insiemi)}

\section{Definzione del problema}
\subsection{Istanza}
\subsubsection{Definizione dell'insieme}
Dato un insieme di attività $X_n = \langle 1, \dots, n \rangle$ con $n \in \mathbb{N}\ n>0$ a cui, per ogni attività dell'insieme sono associati i parametri:
\begin{itemize}
    \item $\mathbf{s_i}$ tempo di inizio dell'attività $i$
    \item $\mathbf{f_i}$ tempo di fine dell'attività $i$
    \item $\mathbf{v_i >0}$ il valore attribuito all'attività $i$
\end{itemize}

\subsubsection{Compatibilità}
Siano $i,j \in X_n$ due attività distinte $(i \not = j)$ diciamo che $i$ e $j$ sono compatibili \textbf{se non si sovrappongono} \uline{$[s_i,f_i) \cap [s_j,f_j) = \emptyset$}

\subsubsection{Funzione di valutazione}
$A \subseteq X_n$ contiene solo attività mutualmente compatibili

Sia $A \subseteq X_n$ qualunque definiamo:

{
    \large
    $$
    V(A) = 
    \begin{cases}
        \sum_{i \in A} v_i & \text{se } A \neq \emptyset \\
        0 & \text{se } A = \emptyset
    \end{cases}
    $$
}

Dato un sottoinsieme, la funzione $V$ gli associa un numero reale (ossia il suo valore)


\begin{itemize}
    \item In caso cui il \uline{sottoinsieme non sia vuoto}, il valore associato sarà la somma di tutti i valori delle sue attività
    \item In caso cui il \uline{sottoinsieme sia vuoto}, il valore associato sarà 0
\end{itemize}

\newpage
\subsubsection{Funzione di Compatibilità}
Dato $A \subseteq X_n$ qualunque definiamo:

{
    \large
    $$
    COMP(A) = 
    \begin{cases}
        true & \text{se } A \not = \emptyset \text{ e } A \text{ contiene attività mutualmente compatibili} \\
        true & \text{se } A = \emptyset \\
        false & \text{altrimenti}
    \end{cases}
    $$
}

Dato un sottoinsieme, la funzione $COMP$ restituisce un valore booleano per definire se l'insieme è composta da elementi mutualmente compatibili
\begin{itemize}
    \item In caso l'insieme non sia vuoto e sia composto da elementi mutualmente compatibili, restituirà il valore $true$
    \item In caso l'insieme sia vuoto, restituirà il valore $true$
    \item In caso l'insieme non sia vuoto e siano presenti al suo interno elementi non mutualmente compatibili, restituirà $false$
\end{itemize}

\subsection{Soluzione}
Stiamo cercando un sottoinsieme di $X_n$ che contenga elementi mutualmente compatibili la cui somma porta al valore massimo

{
    \large
    \[
        V(S) = \max_{\substack{A \subseteq X_n : \\ \text{COMP}(A)=\text{true}}} \{ V(A) \}
    \]
}

\newpage

\section{Definzione della soluzione ricorsiva}

\subsection{Istanza}
Abbiamo un insieme di attività $X_i = \{1, \dots, i\}$


\subsection{Soluzione}
Chiamiamo $S_i$ la \uline{soluzione del sottoproblema i-esimo}

{
    \large
    $$
        S_i \subseteq X_i :  
        V(S_i) = \max_{\substack{A \subseteq X_n : \\ \text{COMP}(A)=\text{true}}} \{ V(A) \}
    $$
}

Vogliamo una soluzione il cui valore sia il massimo dell'insieme degli elementi $A \subseteq X_n$ mutualmente compatibili


\subsubsection{Definizione del coefficiente $c_i$}
Definiamo $c_i$ come $c_i = V(S_i)$

\subsection{Casi base}

\subsubsection{Insieme vuoto $i=0$}
{
    \large
    $$
        i=0 \quad X_0 = \emptyset \implies S_0 = \emptyset \quad c_0=V(S_0)=0
    $$
}

Essendo che l'insieme vuoto non contiene alcuna attività il suo valore sarà pari a 0

\subsubsection{Insieme composto da un solo elemento $i=1$}
{
    \large
    $$
        i=1 \quad X_1=\{1\} \implies S_1 = \{1\} \quad c_1 = V(S_1) = v_1
    $$
}

Essendo l'insieme $S$ composto da un unico elemento esso rispetterà la condizione di compatibilità portando il valore della soluzione ad essere il valore della singola attività presente

\subsection{Passo ricorsivo}
Si vuole risolvere il sottoproblema i-esimo con $i>1$ assumendo di aver già risolto i sottoproblemi $0,\dots,i-1$ ossia di disporre di $S_0,\dots, S_{i-1}$ ma:

{
    \large
    $$
        S_i \overset{???}{=} \underbrace{S_0, \dots, S_{i-1}}_{\text{Espressione in funzione di $S_0,\dots,S_{i-1}$}}
    $$
}
Distinguiamo 2 casi

\subsubsection{$\mathbf{i \not \in S_i}$}
L'elemento i dell'istanza $X_i = \{1,\dots,i\}$ è \textbf{inutile}.

Procediamo come se l'istnaza fosse $X_{i-1}$, quindi procediamo a risolvere il sottoproblema i equivale a risolvere il sottoproblema $i-1$, pertanto:

{
    \large
    $$
        S_i = S_{i-1} \quad c_i = V(S_i)=c_{i-1}
    $$
}

\subsubsection{$\mathbf{i \in S_i}$}
Quindi abbiamo che:

{
    \large
    $$
        S_i = S_? \ \cup \ \{i\}
    $$
}

$S_?$ è la soluzione di quale sottoproblema di $i?$ Per scoprirlo abbiamo bisogno di ordinare per tempo di fine (preprocessing) e calcolare quale delle attività antecendenti ha valore massimo.

Per fare ciò definiamo:

{
    \large

    $$
        p(i) = max\{j<i:\text{$j$ è compatibile con $i$}\}
    $$
}

che ci restituisce l'attività con valore massimo antecedente e compatibile ad $i$.

\begin{itemize}
    \item $\mathbf{X_{p(i)}}$ contiene attività ciascuna delle quali è compatibili con $i$
    \item $\mathbf{S_{p(i)} \subseteq X_{p(i)}}$ contiene attività tutte mutualmente compatibili e ciascuna delle quali è compatibile con $i$, inoltre è di valore massimo
\end{itemize}

Se entrambe queste condizioni vengono rispettate allora:

{
    \large

    $$
        S_i = S_{p(i)} \cup \{i\} \quad c_i = V(S_i) = c_{p(i) + v_i}
    $$
}

\section{Algoritmo bottom-up}
\subsection{Pseudocodice}
\begin{algorithm}
    \caption{WIS-WEIGHT}
    \begin{algorithmic}
        \Procedure{WIS-WEIGHT}{$n$}
            \State $c[0]=0$
            
            \For{$i=1$ to $n$}
                \State $c[i]=max\{c[i-1],c[p(i)] + v_i\}$
            \EndFor

            \State \Return $c$
            
        \EndProcedure
    \end{algorithmic}
\end{algorithm}
\subsection{Spiegazione}
\begin{enumerate}
    \item Impostiamo il caso base $(c[0])$ a 0
    \item Ricorsivamente risaliamo l'array e utilizziamo la funzione di valutazione per assegnare un valore a ognuna delle celle, grazie al fatto che salviamo il valore ottenuto all'interno dell'array $c$ non dovremo ricalcolare ogni volta le soluzioni dei sottoproblemi
\end{enumerate}

\subsection{Complessità}
\begin{itemize}
    \item \textbf{Complessità in tempo:} $\Theta(n)$, è un ciclo for che scorre un array e applica una funzione
    \item \textbf{Complessità in spazio:} $\Theta(n)$ tutto ciò che abbiamo in memoria è un semplice array
\end{itemize}

\subsection{Simulazione}


\newpage
\section{Algoritmo di stampa della soluzione $\mathbf{S_i}$}
\subsection{Pseudocodice}
\begin{algorithm}
    \caption{PRINT-WIS}
    \begin{algorithmic}
        \Procedure{PRINT-WIS}{$i,c$}
            \If{$i=0$}
            \State \Return

            \Else
                \If{$c[p(i)]+v_i \geq c[i-1]$} 
                    \State PRINT-WIS($p(i),c$) \Comment{Stampa $S_{p(i)} \implies S_i=S_{p(i)} \cup$}
                    \State print(i)

                    \Else
                        \State PRINT-WIS($i-1,c$) \Comment{Stampa $S_{i-1}\implies S_i = S_{i-1}$}
                \EndIf
        \EndIf
        \EndProcedure
    \end{algorithmic}
\end{algorithm}
\subsection{Spiegazione}
\begin{enumerate}
    \item Se abbiamo $i=0$ implica che l'insieme è vuoto/siamo arrivati alla fine e quindi non c'è nulla da stampare
    \item Usiamo la formula di ricorrenza per capire quale elemento stampare al prossimo ciclo ricorsivo
    \item Continuiamo finchè $i=0 \implies$ l'array è finito
\end{enumerate}
\subsection{Complessità}
\begin{itemize}
    \item \textbf{Complessità in tempo:} $\Theta(n)$, stiamo solo risalendo l'array ricorsivamente e stampando gli elementi
    \item \textbf{Complessità in spazio:} $\Theta(n)$ tutto ciò che abbiamo in memoria è un semplice array
\end{itemize}


\newpage
\section{Algoritmo di stampa dell'insieme di soluzione}
\subsection{Pseudocodice}
\begin{algorithm}
    \caption{SET-WIS}
    \begin{algorithmic}
        \Procedure{SET-WIS}{$i,c$}
            \If{$i=0$}
                \State \Return $\emptyset$

                \Else
                    \If{$c[p(i)] + v_i \geq c[i-1]$} 
                        \State \Return SET-WIS($p(i),c$) $\cup \{i\}$

                        \Else
                            \State \Return SET-WIS($i-1,c$)
                    \EndIf
            \EndIf
        \EndProcedure
    \end{algorithmic}
\end{algorithm}
\subsection{Spiegazione}
    \begin{enumerate}
        \item Se l'insieme è vuoto allora restituiamo l'insieme vuoto
        \item Attraverso la chiamata formula di ricorsione costruisco l'array contenente gli elementi dell'insieme della soluzione
        \item Finiti gli elementi stampero $\emptyset$ concatenato agli elementi dell'insieme di soluzione
    \end{enumerate}
\subsection{Complessità}
\begin{itemize}
    \item \textbf{Complessità in tempo:} $\Theta(n)$, stiamo solo risalendo l'array ricorsivamente e stampando gli elementi
    \item \textbf{Complessità in spazio:} $\Theta(n)$ tutto ciò che abbiamo in memoria è un semplice array
\end{itemize}

\newpage
\section{Dimostrazione del problema della sottostruttura ottima}
\subsection{Teorema}
\subsection{Dimostrazione}
